\section{模型检验与灵敏度分析}\label{sec:verification}

本节通过解析解对照、空间与时间加密、整体水分收支及静态极限检验数值求解，并利用物性、几何和边界输入的对照分析，评估模型设定对预测结果的影响。各项检验都用未舍入的计算值。本次按最终配置重新完成各问输出场的空间、时间与界面求积对照，以及长时根和独立通量核验；解析级数、向后 Euler 阶数、静态极限、三情形、参数扰动及中截面端效应沿用赛时辅助计算，适用条件逐项列明。数值收敛和离散收支只核验当前数学模型的实现，缺少内部场实测数据时不能据此声称真实预测误差已验证。

\subsection{解析基准与空间收敛性}

\parahead{目的与方法} 在常物性、恒定环境的简化设定下，圆柱的第三类边界瞬态导热有经典的级数解\cite{carslaw}。以其为基准可把空间离散误差单独考察。无量纲温度写成

\begin{equation}
\Theta(\eta,Fo)=\sum_{n} c_n\,J_0(\lambda_n\eta)\,e^{-\lambda_n^2 Fo}
\label{eq:series}
\end{equation}

其中 $\Theta=(T-T_\infty)/(T_0-T_\infty)$ 为无量纲温度、$T_\infty$ 为恒定环境温度，$\eta=r/R_0$ 为无量纲半径，$J_0$、$J_1$ 为第一类零阶、一阶 Bessel 函数，展开系数 $c_n$ 由初始条件正交展开得到。温度辅助问题的 Fourier 数取 $Fo_h=\alpha t/R_0^2$（热扩散率 $\alpha=k/(\rho c_p)$）、Biot 数 $Bi_h=hR_0/k\approx1.39$；含水率辅助问题取常扩散系数 $D\equiv D(C_0)$，其 Fourier 数为 $Fo_m=Dt/R_0^2$、Biot 数 $Bi_m=h_mR_0/D\approx3.24$。两种情形的特征值 $\lambda_n$ 均由 $\lambda_n J_1(\lambda_n)=Bi\,J_0(\lambda_n)$ 确定，$Bi$ 取对应的 $Bi_h$ 或 $Bi_m$。

\parahead{结果} 固定很紧的时间容差、只改变网格时，取 $t=\SI{100}{s}$ 时刻沿半径的最大误差随 $N$ 的变化见图~\ref{fig:aconv} 与表~\ref{tab:aconv}，两者都接近二阶收敛。据此按各问实际输出位置的误差选定计算网格。

\begin{figure}[H]\centering
\includegraphics[width=\textwidth]{fig_analytic_convergence.pdf}
\caption{$t=\SI{100}{s}$ 时数值解相对解析级数解沿半径的最大误差随网格加密的收敛（双对数），虚线为二阶参考斜率；该结论适用于所设常物性、恒环境的解析辅助问题。}
\label{fig:aconv}
\end{figure}

\begin{table}[H]\centering\small
\caption{解析基准误差（$t=\SI{100}{s}$ 时沿半径的最大误差，常物性、恒环境辅助问题）}\label{tab:aconv}
\begin{tabular}{C{2.2cm}C{4.6cm}C{4.6cm}}
\toprule
$N$ & 温度沿半径最大误差 / $\si{\celsius}$ & 含水率沿半径最大误差 / $\si{kg/kg}$ \\
\midrule
200 & $1.0396\times10^{-4}$ & $2.7995\times10^{-4}$ \\
400 & $2.5989\times10^{-5}$ & $6.9857\times10^{-5}$ \\
800 & $6.4991\times10^{-6}$ & $1.7457\times10^{-5}$ \\
\bottomrule
\end{tabular}
\end{table}

\parahead{结论} 解析基准检验显示该常物性辅助问题接近二阶空间收敛，支持对应离散实现的一致性；实际各问的计算网格取值，则根据后文实际输出位置的加密对照确定，烘干时长变化小不能替代完整的场点差异检验。

\subsection{时间离散误差与界面系数}

\parahead{目的与方法} 用向后 Euler 在固定细网格上做时间步加密以检验时间精度；再改变 BDF 的相对容差与界面求积点数，看数值解是否已不随之变化；界面系数则对照调和平均与积分平均两种取法（图~\ref{fig:convergence}）。

\parahead{结果} 图~\ref{fig:bebdf} 给出向后 Euler 对照计算与自适应 BDF 求解在辅助热问题（常物性、$\Tair\equiv\SI{50}{\celsius}$、$N=400$、$0$--$\SI{100}{s}$）上的时间积分对照，参考解为紧容差 BDF（rtol $10^{-11}$、atol $10^{-13}$）。$\SI{100}{s}$ 表面温度相对该参考的差异：向后 Euler 在 $\Delta t=1/0.5/0.25\,\si{s}$ 下依次约为 $8.33\times10^{-3}$、$4.17\times10^{-3}$、$2.08\times10^{-3}\,\si{K}$，\textbf{步长减半后差异近似减半}，观测收敛阶约为 $1.00$；\textbf{自适应 BDF 在当前容差}（rtol $10^{-8}$）\textbf{下与参考结果一致}，差异仅约 $1.29\times10^{-6}\,\si{K}$。本次将 BDF 相对、绝对容差均缩至原值的 $1/100$，并将观测段/延续段最大步长由 $60/600\,\si{s}$ 减至 $30/300\,\si{s}$，问题三、四阈值根分别变化 $\SI{1.758e-8}{h}$、$\SI{5.797e-7}{h}$；独立把界面求积由 $8$ 点增至 $16$ 点，分别变化 $\SI{3.369e-7}{h}$、$\SI{6.241e-7}{h}$。界面系数的对照更为关键：调和平均在 $N=200/400/800$ 给出的烘干时长为 $58.96/57.69/57.51\,\si{h}$，仍随网格明显漂移；积分界面给出 $57.4745/57.4741/57.4740\,\si{h}$，几乎不随网格变化。

\begin{figure}[H]\centering
\includegraphics[width=\textwidth]{fig_be_bdf.pdf}
\caption{辅助热问题的时间积分对照（常物性、恒环境，$N=400$，$0$--$\SI{100}{s}$）：(a) 表面温度时程，各方法在此标度下基本重合；(b) 向后 Euler 的 $\SI{100}{s}$ 终点表面温度相对紧容差 BDF 参考的差异随步长的双对数关系，虚线为一阶参考斜率，自适应 BDF（当前容差）的差异量级单独标示。}
\label{fig:bebdf}
\end{figure}

\begin{figure}[H]\centering
\includegraphics[width=\textwidth]{fig_convergence.pdf}
\caption{烘干时长随网格加密的收敛（问题二、三）：(a) 烘干时长 $t^*$ 随 $N$ 的变化；(b) 各方法烘干时长相对其自身 $N=800$ 结果的变化幅度（对数），表示对网格取值的依赖，而非相对已知真解的误差；积分界面的变化量级远低于调和平均。}
\label{fig:convergence}
\end{figure}

\parahead{结论} 时间容差收紧与界面求积加密引起的烘干时长变化均较小，故本文采用相对容差 $10^{-8}$ 的 BDF、$8$ 点 Gauss 求积。相较调和平均，积分平均在所检验网格上的烘干时长变化更小，因此\textbf{将积分平均用于水分界面扩散系数 $D$ 的计算}；导热系数 $k$ 的界面仍取调和平均。

\subsection{题目指定输出位置的独立加密对照}

\parahead{方法与覆盖} 本次以问题一 $N=1600$、问题二至四 $N=800$、积分界面 $8$ 点、BDF 相对容差 $10^{-8}$ 为基准，分别只加倍空间网格、加强时间控制、加倍界面求积点数，其他设置保持不变。空间对照检查离散和重构共同造成的差异；时间控制与求积对照按上一小节给定设置执行，三类不能互相替代。

问题一逐秒检查 $1$--$\SI{1800}{s}$ 的 $21$ 个径向位置，两场共覆盖 $1800$ 个时刻；问题二、三同样检查两场的 $21$ 个位置，逐秒延伸至 $\SI{206940}{s}$ 并显式并入阈值根，共 $206941$ 个时刻。问题四检查全部逐分钟输出行并并入根时刻，共 $3067$ 个时刻；检查量为域内固定物理位置与移动表面的含水率，域外留空，温度不属于本问结果文件的场验收范围。表中最大差是所有已检查时空场点的最大绝对差，不是两个最大含水率标量之差。

\begin{table}[H]\centering\small
\caption{本次复算的场点与阈值根加密差异}\label{tab:gridsum}
\begin{tabular}{C{1.3cm}C{3.0cm}C{3.0cm}C{2.9cm}C{2.9cm}}
\toprule
问题 & 独立对照 & 最大含水率差 / $\si{kg/kg}$ & 最大温度差 / $\si{\celsius}$ & 阈值根差 / $\si{h}$ \\
\midrule
一 & $N:1600\to3200$ & $4.693\times10^{-5}$ & $2.232\times10^{-4}$ & \textemdash \\
一 & 时间控制加密 & $2.473\times10^{-7}$ & $2.279\times10^{-4}$ & \textemdash \\
一 & 求积 $8\to16$ 点 & $2.045\times10^{-13}$ & $5.645\times10^{-10}$ & \textemdash \\
二、三 & $N:800\to1600$ & $1.549\times10^{-4}$ & $2.012\times10^{-4}$ & $2.219\times10^{-5}$ \\
二、三 & 时间控制加密 & $4.959\times10^{-7}$ & $2.230\times10^{-4}$ & $1.758\times10^{-8}$ \\
二、三 & 求积 $8\to16$ 点 & $5.406\times10^{-7}$ & $2.083\times10^{-4}$ & $3.369\times10^{-7}$ \\
四 & $N:800\to1600$ & $2.194\times10^{-4}$ & \textemdash & $2.026\times10^{-5}$ \\
四 & 时间控制加密 & $6.521\times10^{-7}$ & \textemdash & $5.797\times10^{-7}$ \\
四 & 求积 $8\to16$ 点 & $5.561\times10^{-7}$ & \textemdash & $6.241\times10^{-7}$ \\
\bottomrule
\end{tabular}
\end{table}

\parahead{结果与解释} 所有适用场点差均低于预设比较尺度：含水率 $\SI{5e-4}{kg/kg}$、温度 $\SI{1e-2}{\celsius}$；根时刻差均小于 $\SI{0.02}{h}$。问题四空间差最大处为 $t=\SI{124620}{s}$、$r=\SI{1.2}{cm}$。这些结果支持采用上述计算配置，但表中的差异是两次离散解之差，不是相对未知真解的认证误差；输出四位小数也不表示四位数字均已获得物理准确性验证。

\subsection{守恒与独立通量核算}

\parahead{目的与方法} 与模型水分收支直接对应的核验，是把表面失水速率沿时间独立求积，与平均干基含水率的实际下降量对照（离散格式的代数自洽性另在附录核算）。平均干基含水率在固定域取 $\Cbar_{23}=\tfrac{2}{R_0^2}\sum_iV_iC_i$；问题四在参考坐标下取

\begin{equation}
W_i=\int_{\Omega_i^x}x\,\mathrm{d}x,\qquad \sum_iW_i=\frac12,\qquad \Cbar_4(t)=2\sum_iW_i\,\tilde c_i(t)
\label{eq:cbar4}
\end{equation}

其中 $\Omega_i^x$ 为参考坐标中的控制区间，$W_i$ 为对应的几何权重；在干物质基准密度径向均匀的假设下，$\Cbar_4$ 也是干物质质量加权的平均含水率。由式~\eqref{eq:massbal}，平均干基含水率的损失应等于表面通量的累计。连续解的独立求积在固定域用初始半径 $R_0$ 与表面值 $C(R_0,\tau)$：

\begin{equation}
I_{23}(t)=\int_0^{t}\frac{2h_m}{R_0}\big[C(R_0,\tau)-\Cenv(\tau)\big]\,\mathrm{d}\tau
\label{eq:flux23}
\end{equation}

问题四则用当前半径 $R(\tau)$ 与参考坐标下的表面值 $\tilde c(1,\tau)$：

\begin{equation}
I_{4}(t)=\int_0^{t}\frac{2h_m}{R(\tau)}\big[\tilde c(1,\tau)-\Cenv(\tau)\big]\,\mathrm{d}\tau
\label{eq:flux4}
\end{equation}

\parahead{结果} 将表面失水速率换算为单位干物质质量的失水速率，再沿时间积分，得到平均干基含水率的累计损失，并与内部平均含水率的实际下降量比较（连续解的表面通量按式~\eqref{eq:flux23}、\eqref{eq:flux4} 独立求积）。本次从初值独立求积至各问阈值根，计算网格 $N=800$ 下问题二、三与问题四的相对差分别为 \textbf{$1.058\times10^{-8}$} 和 \textbf{$1.107\times10^{-8}$}；加密至 $N=1600$ 后分别为 $1.131\times10^{-8}$、$9.115\times10^{-9}$，均小于 $10^{-7}$，表明内部含水率的减少与表面累计失水相对应，数值解维持了所建模型的整体水分收支。赛时辅助计算的离散代数自洽核算（向后 Euler 逐步收支约 $\num{1.5e-14}$、离散热收支残差约 $\SI{3.212e-10}{W/m}$）见附录~\ref{sec:divexp}。

\subsection{物理界限与静态极限}

\parahead{目的与方法} 赛时辅助包络检验取 $N=200$，问题二、三的时间窗为 $0$--$\SI{8000}{s}$，问题四为 $0$--$\SI{3600}{s}$，各取 $25$ 个采样时刻，将重构场与初边值历史包络比较。静态极限检验同取 $N=200$、积分界面，在附录 4 物性与相同初边值下令 $R\equiv R_0$，于 $\SI{3600}{s}$ 比较动域和固定域求解器。

\parahead{结果与结论} 上述时间窗的 $25$ 个采样时刻均未发现超过设定容差的包络越界，这不等于验证全部接受时间层或完整干燥过程。静态极限对照的相对差约 $\num{3.6e-8}$，支持式~\eqref{eq:q4govc}--\eqref{eq:q4govt} 在该配置与时刻退化到固定域形式。

\subsection{物性与几何效应的分离}

\parahead{目的与方法} 问题四相对问题三同时改变了物性（附录 3$\to$4）和几何（固定$\to$收缩），需把两者分开。为此增加「附录 4 物性、半径固定为 $R_0$」的对照，与另两种情形在同一网格上比较（图~\ref{fig:threecase}）。

\begin{figure}[H]\centering
\includegraphics[width=0.82\textwidth]{fig_threecase.pdf}
\caption{三种情形的烘干时长（同一网格 $N=400$）：(1)附录3 物性、固定 $R_0$；(2)附录4 物性、固定 $R_0$；(3)附录4 物性、收缩 $R(t)$。物性变化（(1)→(2)）延长烘干，尺寸收缩（(2)→(3)）缩短烘干。}
\label{fig:threecase}
\end{figure}

\parahead{结果与结论} 三种情形的烘干时长分别约为 \textbf{$\SI{57.4741}{h}$}、\textbf{$\SI{129.8487}{h}$}、\textbf{$\SI{51.0921}{h}$}（同网格 $N=400$）。固定半径时把物性从附录 3 换成附录 4，烘干时长增加约 $\SI{72.37}{h}$；在附录 4 物性下再加入给定收缩，时长减少约 $\SI{78.76}{h}$，两步合计使问题三到问题四缩短约 $\SI{6.38}{h}$。物性变化与几何收缩沿这一比较路径的作用方向相反；此为既定比较条件下的分解，并不宣称两种作用彼此独立。这三个值属于同网格 $N=400$ 的辅助对照，与两问最终 $N=800$ 的结果分开说明。

\subsection{端面影响的量级估计}

\parahead{目的与方法} 该估计在常物性、恒环境、附录 2 热参数与固定 $D(C_0)$ 的线性辅助问题上进行，采样位置为固定圆柱的中截面。用乘积解 $\Theta_{\mathrm{finite}}=\Theta_{\mathrm{radial}}\,\Theta_{\mathrm{slab}}$ 近似有限长圆柱，端面带来的温度增量为 $\Delta T=(T_0-T_\infty)\,\Theta_{\mathrm{radial}}(\Theta_{\mathrm{slab}}-1)$（含水率同理），其中 $\Theta_{\mathrm{radial}}$、$\Theta_{\mathrm{slab}}$ 分别为径向与平板方向的无量纲解，在中截面各相关时刻与位置上取最大。该对照使用附录 2 热参数，不代表已求解四问各自的轴向变物性模型。

\parahead{结果与结论} 在所设对照问题的中截面采样点，计入两端换热与传质后，温度和含水率的最大绝对变化分别为 \textbf{$\SI{9.10e-4}{\celsius}$}（$\SI{1.5}{h}$、中心）和 \textbf{$\num{1.13e-5}\,\si{kg/kg}$}（$\SI{42}{h}$、中心），均低于所设比较尺度（$\SI{0.1}{\celsius}$、$0.005$）；含水率增量在平板本征项数取 $120$ 与 $200$ 时稳定于 $\num{1.13e-5}$。机理上，轴向扩散时间尺度约 $(L/2)^2/D\approx\SI{886}{h}$，远大于径向的 $R_0^2/D\approx\SI{22.7}{h}$，端面效应显现时径向场早已接近均衡。该结果仅表明所设线性辅助问题的中截面采样位置端效应较弱，为一维近似提供量级参考；它不覆盖端部，也不构成四问变物性全圆柱的误差界。

\subsection{参数扰动分析}

\parahead{目的与方法} 以问题二、三为基线，逐一对参数作情景扰动并重新求解，观察烘干时长的变化（图~\ref{fig:sensitivity}、表~\ref{tab:sens}）。三类情景的幅度与来源见表~\ref{tab:sensrc}：换热系数 $h$、传质系数 $h_m$ 取题给值的 $0.5$ 倍与 $2$ 倍；外推空气温度与环境含水浓度的扰动幅度，由附件 1 中 $2$--$\SI{4}{h}$ 平台数据相对末小时均值的最大偏差确定，分别为 $\pm\SI{0.39}{\celsius}$ 与 $\pm0.0011\,\si{kg/kg}$，仅施加于 $\SI{4}{h}$ 之后的外推段；另用末值延续与局部平滑考察输入处理方式。基线烘干时长为 $\SI{57.4741}{h}$（$N=400$）。图中淡灰带 $\pm\SI{0.02}{h}$（约 $1.2\,\mathrm{min}$）是取自验收配置的统一时长变化筛选尺度，用以突出主要变化，并非数值误差量级。

\begin{table}[H]\centering\small
\caption{参数扰动情景的幅度与来源}\label{tab:sensrc}
\begin{tabular}{C{3.0cm}C{7.6cm}C{2.4cm}}
\toprule
情景 & 幅度与来源 & 作用范围 \\
\midrule
换热系数 $h$ & 题给 $\SI{25}{W/(m^2.K)}$ 的 $0.5$、$2$ 倍 & 全过程 \\
传质系数 $h_m$ & 题给 $\SI{8e-7}{m/s}$ 的 $0.5$、$2$ 倍 & 全过程 \\
外推空气温度 & $2$--$\SI{4}{h}$ 平台相对末小时均值的最大偏差 $\pm\SI{0.39}{\celsius}$ & $\SI{4}{h}$ 之后 \\
外推环境含水浓度 & 同法得 $\pm0.0011\,\si{kg/kg}$ & $\SI{4}{h}$ 之后 \\
末值延续 & 以末观测点代替末小时均值 & $\SI{4}{h}$ 之后 \\
局部平滑 & 内点 $(y_{-}{+}2y{+}y_{+})/4$ 后再线性插值 & $0$--$\SI{4}{h}$ \\
\bottomrule
\end{tabular}
\end{table}

\begin{figure}[H]\centering
\includegraphics[width=0.96\textwidth]{fig_sensitivity.pdf}
\caption{参数扰动对烘干时长的影响：(a) 全部情景；(b) 小影响情景放大，淡灰带为 $\pm0.02$ h 的时长变化筛选尺度。}
\label{fig:sensitivity}
\end{figure}

\begin{table}[H]\centering\small
\caption{参数扰动的烘干时长变化}\label{tab:sens}
\begin{tabular}{C{6.0cm}C{3.0cm}C{3.0cm}}
\toprule
扰动情景 & 时长 / $\si{h}$ & 时长变化 / $\si{h}$ \\
\midrule
换热系数 $h$ 减半 & 57.5537 & $+0.0796$ \\
换热系数 $h$ 加倍 & 57.4354 & $-0.0387$ \\
传质系数 $h_m$ 减半 & 64.9292 & $+7.4550$ \\
传质系数 $h_m$ 加倍 & 54.6719 & $-2.8022$ \\
外推空气温度 $+0.39$ °C & 56.7682 & $-0.7059$ \\
外推空气温度 $-0.39$ °C & 58.1924 & $+0.7183$ \\
外推环境含水浓度 $+0.0011$ & 57.4965 & $+0.0224$ \\
外推环境含水浓度 $-0.0011$ & 57.4526 & $-0.0215$ \\
外推改为末值保持 & 57.1695 & $-0.3046$ \\
数据平滑 smooth121 & 57.4741 & $+0.0000$ \\
\bottomrule
\end{tabular}
\end{table}

\parahead{结果与结论} 在问题二、三的上述情景中，应优先核实表面传质系数，并补充 $\SI{4}{h}$ 之后的温度观测。\textbf{传质系数减半使烘干时长增加约 $13.0\%$、加倍则缩短约 $4.88\%$}，两种等倍率变化的作用并不对称；换热系数同样倍率变化引起的时长差在 $\SI{0.1}{h}$ 以内。外推温度变化 $\pm\SI{0.39}{\celsius}$ 引起约 $42$--$43\,\mathrm{min}$ 的时长变化、末值延续约 $18.3\,\mathrm{min}$，环境含水浓度与局部平滑更小。因而，提高传质参数与后期环境输入的可信度，比进一步平滑现有分钟级数据更有助于控制本组预测中的输入不确定性。这些情景是对既定输入假设的条件比较，不是时长置信区间；$h$ 的小幅时长响应也不能验证未单列潜热的热边界。
